\begin{table}[ht]
\centering
\caption{News Media Consumption}
\label{tab:followup_media}
\vspace{6pt}
\newlength{\mytabwidth}
\settowidth{\mytabwidth}{%
\begin{tabular}{l >{\centering\arraybackslash}p{3cm} >{\centering\arraybackslash}p{3.5cm}}
\toprule
\textbf{Media Source} & \textbf{Respondents} & \textbf{\% of Follow-Up Sample} \\
\bottomrule
\end{tabular}}
\begin{tabular}{l >{\centering\arraybackslash}p{3cm} >{\centering\arraybackslash}p{3.5cm}}
\toprule
\textbf{Media Source} & \textbf{Respondents} & \textbf{\% of Sample} \\
\midrule
Süddeutsche Zeitung & 1,661 & 97.1 \\
ARD/ZDF & 1,475 & 86.3 \\
Die Zeit & 520 & 30.4 \\
Der Spiegel & 503 & 29.4 \\
FAZ & 239 & 14.0 \\
TAZ & 225 & 13.2 \\
NZZ & 143 & 8.4 \\
Handelsblatt & 134 & 7.8 \\
Instagram & 121 & 7.1 \\
YouTube & 109 & 6.4 \\
Bild & 54 & 3.2 \\
X (Twitter) & 20 & 1.2 \\
TikTok & 6 & 0.4 \\
\bottomrule
\end{tabular}

\vspace{4pt}
\parbox{\mytabwidth}{\scriptsize \textit{Note:} The survey question in the follow-up survey from February 2026 was:  An individual could opt for as many outlets from the above list as they wanted. N = 1,710 follow-up respondents.}
\end{table}
